% ==============================================================================
% BỘ ĐỀ THAM KHẢO TUYỂN SINH 10 NĂM HỌC 2026-2027 SỞ GD&ĐT TP.HCM
% ĐỢT 12: CỤM 8 - ĐỀ 3 & CỤM 9 - ĐỀ 1 (14 CÂU TỰ LUẬN)
% ==============================================================================

% ==============================================================================
% CỤM 8 - ĐỀ THAM KHẢO 3 (Đề số 24)
% ==============================================================================

\begin{bt}%[9D6H18-2]
% nguon: Đề tham khảo Tuyển sinh 10 năm học 2026-2027 Sở GD&ĐT TP.HCM - Cụm 8 Đề 3, Bài 1
	Cho parabol $(P)\colon y = \dfrac{1}{3}x^2$.
	\begin{enumerate}[a)]
		\item Vẽ đồ thị $(P)$ trên hệ trục tọa độ $Oxy$.
		\item Tìm tọa độ các điểm thuộc $(P)$ có hoành độ và tung độ là hai số đối nhau.
	\end{enumerate}
	\loigiai{
		\begin{enumerate}[a)]
			\item Bảng giá trị tương ứng giữa $x$ và $y$ của parabol $(P)\colon y = \dfrac{1}{3}x^2$:
			\begin{center}
				\begin{tabular}{|c|c|c|c|c|c|}
					\hline
					$x$ & $-6$ & $-3$ & $0$ & $3$ & $6$ \\
					\hline
					$y = \dfrac{1}{3}x^2$ & $12$ & $3$ & $0$ & $3$ & $12$ \\
					\hline
				\end{tabular}
			\end{center}
			Đồ thị $(P)$ là một parabol đi qua gốc tọa độ $O(0;0)$, nằm phía trên trục hoành và nhận trục tung $Oy$ làm trục đối xứng.
			\begin{center}
				\begin{tikzpicture}[>=stealth, scale=0.6]
					\draw[->] (-7,0) -- (7,0) node[below] {$x$};
					\draw[->] (0,-1) -- (0,13) node[left] {$y$};
					\fill (0,0) circle (1.5pt) node[below right] {$O$};
					\draw[smooth, samples=100, domain=-6.1:6.1, thick, color=blue] plot(\x, {1/3*(\x)^2});
					\draw[dashed] (-6,0) node[below] {$-6$} -- (-6,12) -- (0,12) node[left] {$12$} -- (6,12) -- (6,0) node[below] {$6$};
					\draw[dashed] (-3,0) node[below] {$-3$} -- (-3,3) -- (0,3) node[left] {$3$} -- (3,3) -- (3,0) node[below] {$3$};
					\fill (-6,12) circle (1.5pt) (6,12) circle (1.5pt) (-3,3) circle (1.5pt) (3,3) circle (1.5pt);
				\end{tikzpicture}
			\end{center}
			\item Vì hoành độ và tung độ là hai số đối nhau nên ta có phương trình:
			$$y = -x.$$
			Mặt khác điểm đó thuộc $(P)$ nên:
			$$\dfrac{1}{3}x^2 = -x \iff \dfrac{1}{3}x^2 + x = 0 \iff x\left(\dfrac{1}{3}x + 1\right) = 0 \iff \left[\begin{aligned} &x = 0 \\ &x = -3. \end{aligned}\right.$$
			- Với $x = 0 \implies y = 0$.
			
			- Với $x = -3 \implies y = 3$.
			
			Vậy tọa độ các điểm cần tìm là $(0; 0)$ và $(-3; 3)$.
		\end{enumerate}
	}
\end{bt}

\begin{bt}%[9D6H20-2]
% nguon: Đề tham khảo Tuyển sinh 10 năm học 2026-2027 Sở GD&ĐT TP.HCM - Cụm 8 Đề 3, Bài 2
	Cho phương trình $-5x^2 + 3x + 4 = 0$.
	\begin{enumerate}[a)]
		\item Chứng minh phương trình trên có hai nghiệm phân biệt $x_1, x_2$.
		\item Không giải phương trình, hãy tính giá trị của biểu thức:
		$$A = \dfrac{2x_1}{x_2} + \dfrac{2x_2}{x_1}.$$
	\end{enumerate}
	\loigiai{
		\begin{enumerate}[a)]
			\item Phương trình $-5x^2 + 3x + 4 = 0$ có $a = -5, b = 3, c = 4$.
			
			Vì $a \cdot c = (-5) \cdot 4 = -20 < 0$ nên phương trình luôn có hai nghiệm phân biệt $x_1, x_2$.
			\item Theo định lí Viète, ta có:
			$$\begin{cases} x_1 + x_2 = \dfrac{3}{5} \\ x_1 x_2 = -\dfrac{4}{5}. \end{cases}$$
			Ta có:
			$$x_1^2 + x_2^2 = (x_1 + x_2)^2 - 2x_1 x_2 = \left(\dfrac{3}{5}\right)^2 - 2 \cdot \left(-\dfrac{4}{5}\right) = \dfrac{9}{25} + \dfrac{8}{5} = \dfrac{49}{25}.$$
			Quy đồng mẫu số và biến đổi biểu thức $A$:
			$$A = \dfrac{2x_1^2 + 2x_2^2}{x_1 x_2} = \dfrac{2(x_1^2 + x_2^2)}{x_1 x_2}.$$
			Thay các giá trị vào biểu thức $A$, ta được:
			$$A = \dfrac{2 \cdot \dfrac{49}{25}}{-\dfrac{4}{5}} = \dfrac{\dfrac{98}{25}}{-\dfrac{4}{5}} = -\dfrac{98}{25} \cdot \dfrac{5}{4} = -\dfrac{49}{10} = -4{,}9.$$
			Vậy $A = -\dfrac{49}{10}$.
		\end{enumerate}
	}
\end{bt}

\begin{bt}%[9D8H26-3]
% nguon: Đề tham khảo Tuyển sinh 10 năm học 2026-2027 Sở GD&ĐT TP.HCM - Cụm 8 Đề 3, Bài 3
	Trong một lớp có $40$ học sinh, kết quả kiểm tra môn Toán được phân loại như sau: Loại Giỏi có $8$ học sinh; Loại Khá có $14$ học sinh; Loại Đạt có $12$ học sinh; Loại Chưa đạt có $6$ học sinh.
	\begin{enumerate}[a)]
		\item Tính tỉ lệ phần trăm học sinh ở mỗi loại học lực.
		\item Chọn ngẫu nhiên $1$ học sinh trong lớp. Tính xác suất của biến cố: ``Học sinh được chọn thuộc loại Khá hoặc Giỏi''.
		\item Chọn ngẫu nhiên $1$ học sinh trong lớp. Biết rằng học sinh được chọn không thuộc loại Chưa đạt, tính xác suất của biến cố: ``Học sinh đó thuộc loại Giỏi''.
	\end{enumerate}
	\loigiai{
		\begin{enumerate}[a)]
			\item Tỉ lệ phần trăm học sinh ở từng loại học lực:
			- Loại Giỏi: $\dfrac{8}{40} \cdot 100\% = 20\%$;
			- Loại Khá: $\dfrac{14}{40} \cdot 100\% = 35\%$;
			- Loại Đạt: $\dfrac{12}{40} \cdot 100\% = 30\%$;
			- Loại Chưa đạt: $\dfrac{6}{40} \cdot 100\% = 15\%$.
			\item Số học sinh trong lớp là $n(\Omega) = 40$.
			
			Số học sinh đạt loại Khá hoặc Giỏi là:
			$$n(A) = 14 + 8 = 22\text{ (học sinh)}.$$
			Xác suất của biến cố $A$ là:
			$$P(A) = \dfrac{n(A)}{n(\Omega)} = \dfrac{22}{40} = \dfrac{11}{20} = 0{,}55.$$
			\item Vì học sinh được chọn không thuộc loại Chưa đạt nên số kết quả có thể xảy ra là:
			$$n(\Omega') = 40 - 6 = 34\text{ (học sinh)}.$$
			Số học sinh thuộc loại Giỏi là $8$ học sinh, nên số kết quả thuận lợi là $n(B) = 8$.
			
			Xác suất của biến cố $B$ là:
			$$P(B) = \dfrac{n(B)}{n(\Omega')} = \dfrac{8}{34} = \dfrac{4}{17}.$$
		\end{enumerate}
	}
\end{bt}

\begin{bt}%[9D6V21-3]
% nguon: Đề tham khảo Tuyển sinh 10 năm học 2026-2027 Sở GD&ĐT TP.HCM - Cụm 8 Đề 3, Bài 4
	Một khu đất hình chữ nhật có kích thước chiều dài là $x + 7\text{ (m)}$ và chiều rộng là $x + 3\text{ (m)}$. Người ta dành một phần đất hình vuông cạnh $x\text{ (m)}$ ($x > 0$) để trồng cây, phần diện tích còn lại dùng để lát gạch.
	\begin{enumerate}[a)]
		\item Lập biểu thức thu gọn $A$ theo $x$ tính diện tích phần đất để lát gạch.
		\item Biết diện tích phần đất lát gạch là $51\text{ m}^2$, hãy tính chu vi khu đất trồng cây.
	\end{enumerate}
	\loigiai{
		\begin{enumerate}[a)]
			\item Diện tích toàn bộ khu đất hình chữ nhật là:
			$$S_{\text{khu đất}} = (x + 7)(x + 3) = x^2 + 10x + 21\text{ (m}^2).$$
			Diện tích phần đất hình vuông trồng cây là:
			$$S_{\text{cây}} = x^2\text{ (m}^2).$$
			Biểu thức $A$ tính diện tích phần đất lát gạch là:
			$$A = S_{\text{khu đất}} - S_{\text{cây}} = (x^2 + 10x + 21) - x^2 = 10x + 21\text{ (m}^2).$$
			\item Theo đề bài, diện tích phần đất lát gạch là $51\text{ m}^2$ nên ta có phương trình:
			$$10x + 21 = 51 \iff 10x = 30 \iff x = 3\text{ (thỏa mãn } x > 0).$$
			Độ dài cạnh khu đất hình vuông trồng cây là $x = 3\text{ m}$.
			
			Chu vi của khu đất trồng cây là:
			$$C = 4x = 4 \cdot 3 = 12\text{ (m)}.$$
		\end{enumerate}
	}
\end{bt}

\begin{bt}%[9H0V31-3]
% nguon: Đề tham khảo Tuyển sinh 10 năm học 2026-2027 Sở GD&ĐT TP.HCM - Cụm 8 Đề 3, Bài 5
	Một bồn chứa nước gồm hai phần ghép liền nhau: Phần dưới là hình trụ có bán kính đáy $r = 1{,}5\text{ m}$ và chiều cao $h\text{ (m)}$. Phần trên là hình nón có chung đáy với hình trụ và chiều cao là $\dfrac{h}{2}\text{ (m)}$.
	\begin{enumerate}[a)]
		\item Viết biểu thức tính thể tích bồn chứa $V$ theo $h$.
		\item Biết thể tích bồn chứa là $V = 9\pi\text{ m}^3$. Người ta sơn toàn bộ mặt ngoài của bồn chứa (không sơn mặt đáy dưới). Tính diện tích cần sơn của bồn chứa (kết quả làm tròn đến hàng đơn vị của $\text{m}^2$).
	\end{enumerate}
	\loigiai{
		\begin{enumerate}[a)]
			\item Thể tích phần hình trụ là:
			$$V_{\text{trụ}} = \pi r^2 h = \pi \cdot (1{,}5)^2 \cdot h = \dfrac{9}{4}\pi h\text{ (m}^3).$$
			Thể tích phần hình nón là:
			$$V_{\text{nón}} = \dfrac{1}{3}\pi r^2 \cdot \dfrac{h}{2} = \dfrac{1}{3}\pi \cdot (1{,}5)^2 \cdot \dfrac{h}{2} = \dfrac{3}{8}\pi h\text{ (m}^3).$$
			Biểu thức thể tích bồn chứa $V$ theo $h$ là:
			$$V = V_{\text{trụ}} + V_{\text{nón}} = \dfrac{9}{4}\pi h + \dfrac{3}{8}\pi h = \dfrac{21}{8}\pi h\text{ (m}^3).$$
			\item Theo đề bài, ta có:
			$$\dfrac{21}{8}\pi h = 9\pi \iff h = \dfrac{9 \cdot 8}{21} = \dfrac{24}{7}\text{ (m)}.$$
			Diện tích xung quanh của phần hình trụ là:
			$$S_{\text{xq trụ}} = 2\pi r h = 2\pi \cdot 1{,}5 \cdot \dfrac{24}{7} = \dfrac{72}{7}\pi\text{ (m}^2).$$
			Chiều cao phần hình nón là:
			$$h_{\text{nón}} = \dfrac{h}{2} = \dfrac{12}{7}\text{ (m)}.$$
			Đường sinh của phần hình nón là:
			$$l = \sqrt{r^2 + h_{\text{nón}}^2} = \sqrt{(1{,}5)^2 + \left(\dfrac{12}{7}\right)^2} = \sqrt{\dfrac{9}{4} + \dfrac{144}{49}} = \sqrt{\dfrac{1017}{196}} = \dfrac{3\sqrt{113}}{14}\text{ (m)}.$$
			Diện tích xung quanh phần hình nón là:
			$$S_{\text{xq nón}} = \pi r l = \pi \cdot 1{,}5 \cdot \dfrac{3\sqrt{113}}{14} = \dfrac{9\sqrt{113}}{28}\pi\text{ (m}^2).$$
			Diện tích cần sơn của bồn chứa là:
			$$S = S_{\text{xq trụ}} + S_{\text{xq nón}} = \dfrac{72}{7}\pi + \dfrac{9\sqrt{113}}{28}\pi \approx 32{,}317 + 10{,}730 = 43{,}047 \approx 43\text{ (m}^2).$$
			Vậy diện tích cần sơn khoảng $43\text{ m}^2$.
		\end{enumerate}
	}
\end{bt}

\begin{bt}%[9D1H3-4]
% nguon: Đề tham khảo Tuyển sinh 10 năm học 2026-2027 Sở GD&ĐT TP.HCM - Cụm 8 Đề 3, Bài 6
	Một đoàn tàu hỏa chạy với vận tốc không đổi đi qua một cây cầu dài $600\text{ m}$. Từ lúc đầu tàu bắt đầu lên cầu cho đến khi toa cuối cùng ra khỏi cầu mất đúng $40$ giây. Cũng đoàn tàu đó, khi vượt qua một người lái xe mô tô đang chạy cùng chiều với vận tốc $36\text{ km/h}$ mất $20$ giây (tính từ lúc đầu tàu ngang đuôi xe đến khi toa cuối cùng vượt qua đầu xe). Coi xe mô tô là một chất điểm. Tính vận tốc và chiều dài của đoàn tàu.
	\loigiai{
		Gọi $x$ (m/s) là vận tốc của đoàn tàu ($x > 10$).
		
		Gọi $y$ (m) là chiều dài của đoàn tàu ($y > 0$).
		
		Đổi $36\text{ km/h} = \dfrac{36 \cdot 1000}{3600} = 10\text{ m/s}$.
		
		Khi tàu đi qua cầu dài $600\text{ m}$, quãng đường đầu tàu đi được là $y + 600\text{ m}$.
		
		Vì thời gian qua cầu là $40$ giây nên ta có phương trình:
		$$40x = y + 600 \iff 40x - y = 600. \quad (1)$$
		Khi tàu vượt qua xe mô tô chạy cùng chiều với vận tốc $10\text{ m/s}$, trong $20$ giây xe mô tô đi được:
		$$20 \cdot 10 = 200\text{ (m)}.$$
		Quãng đường tàu đi được trong $20$ giây bằng quãng đường mô tô đi được cộng với chiều dài đoàn tàu:
		$$20x = 200 + y \iff 20x - y = 200. \quad (2)$$
		Từ $(1)$ và $(2)$, ta có hệ phương trình:
		$$\begin{cases} 40x - y = 600 \\ 20x - y = 200 \end{cases} \iff \begin{cases} 20x = 400 \\ y = 20x - 200 \end{cases} \iff \begin{cases} x = 20\text{ (thỏa mãn)} \\ y = 200\text{ (thỏa mãn)}. \end{cases}$$
		Vậy vận tốc của đoàn tàu là $20\text{ m/s}$ (tương đương $72\text{ km/h}$) và chiều dài của đoàn tàu là $200\text{ m}$.
	}
\end{bt}

\begin{bt}%[9H9C29-3]
% nguon: Đề tham khảo Tuyển sinh 10 năm học 2026-2027 Sở GD&ĐT TP.HCM - Cụm 8 Đề 3, Bài 7
	Cho đường tròn $(O; R)$ và điểm $A$ nằm ngoài đường tròn sao cho $OA > 2R$. Kẻ hai tiếp tuyến $AB, AC$ với $(O)$ ($B, C$ là các tiếp điểm). Đoạn thẳng $OA$ cắt dây cung $BC$ tại $H$. Vẽ đường kính $CD$ của đường tròn $(O)$. Đường thẳng $AD$ cắt $(O)$ tại điểm thứ hai là $E$ ($E$ khác $D$).
	\begin{enumerate}[a)]
		\item Chứng minh bốn điểm $A, B, O, C$ cùng thuộc một đường tròn và $BD \parallel OA$.
		\item Chứng minh hệ thức $AH \cdot AO = AE \cdot AD$ và $\widehat{AHE} = \widehat{ADO}$.
		\item Cho $OA = 3R$. Tính theo $R$ độ dài đoạn thẳng $BC$ và diện tích tứ giác $ABOC$.
	\end{enumerate}
	\loigiai{
		\begin{center}
			\begin{tikzpicture}[>=stealth, scale=0.9]
				\coordinate (O) at (0,0);
				\def\r{2.2}
				\draw (O) circle (\r);
				\coordinate (A) at (6.6,0);
				\pgfmathsetmacro{\distA}{6.6}
				\pgfmathsetmacro{\alphaA}{acos(\r/\distA)}
				\coordinate (B) at ({\alphaA}:\r);
				\coordinate (C) at ({-\alphaA}:\r);
				\draw (A) -- (B);
				\draw (A) -- (C);
				\draw (O) -- (B);
				\draw (O) -- (C);
				\draw (O) -- (A);
				\draw (B) -- (C);
				\coordinate (H) at (intersection of O--A and B--C);
				\coordinate (D) at ($(C)!2!(O)$);
				\draw (C) -- (D);
				\draw (B) -- (D);
				\draw (A) -- (D);
				\coordinate (E) at ($(A)!0.4!(D)$);
				\draw (H) -- (E);
				\node at (O) [below left] {$O$};
				\node at (A) [right] {$A$};
				\node at (B) [above] {$B$};
				\node at (C) [below] {$C$};
				\node at (D) [above left] {$D$};
				\node at (H) [below right] {$H$};
				\node at (E) [above] {$E$};
				\fill (O) circle (1.5pt) (A) circle (1.5pt) (B) circle (1.5pt) (C) circle (1.5pt) (D) circle (1.5pt) (H) circle (1.5pt);
			\end{tikzpicture}
		\end{center}
		\begin{enumerate}[a)]
			\item Vì $AB, AC$ là hai tiếp tuyến của $(O)$ nên $\widehat{ABO} = 90^\circ$ và $\widehat{ACO} = 90^\circ$.
			
			Tứ giác $ABOC$ có $\widehat{ABO} + \widehat{ACO} = 90^\circ + 90^\circ = 180^\circ$, suy ra bốn điểm $A, B, O, C$ cùng thuộc đường tròn đường kính $AO$.
			
			Vì $AB = AC$ và $OB = OC = R$ nên $OA$ là đường trung trực của $BC \implies OA \perp BC$ tại $H$.
			
			Mặt khác $\widehat{DBC} = 90^\circ$ (góc nội tiếp chắn nửa đường tròn đường kính $CD$) $\implies BD \perp BC$.
			
			Từ đó suy ra $BD \parallel OA$ (vì cùng vuông góc với $BC$).
			\item Xét tam giác vuông $ABO$ tại $B$ có đường cao $BH$:
			$$AB^2 = AH \cdot AO. \quad (1)$$
			Xét $\Delta ACE$ và $\Delta ADC$ có $\widehat{CAD}$ chung và $\widehat{AEC} = \widehat{ACD} = 90^\circ$ (do $\widehat{CED} = 90^\circ$).
			
			Do đó $\Delta ACE \backsim \Delta ADC\text{ (g-g)}$, suy ra:
			$$\dfrac{AC}{AD} = \dfrac{AE}{AC} \implies AC^2 = AE \cdot AD. \quad (2)$$
			Vì $AB = AC$ nên từ $(1)$ và $(2)$ suy ra:
			$$AH \cdot AO = AE \cdot AD.$$
			Từ $AH \cdot AO = AE \cdot AD \implies \dfrac{AH}{AD} = \dfrac{AE}{AO}$.
			
			Xét $\Delta AHE$ và $\Delta ADO$ có:
			$$\widehat{OAD}\text{ chung},\quad \dfrac{AH}{AD} = \dfrac{AE}{AO}.$$
			Do đó $\Delta AHE \backsim \Delta ADO\text{ (c-g-c)} \implies \widehat{AHE} = \widehat{ADO}$.
			\item Khi $OA = 3R$:
			
			Trong tam giác vuông $ABO$ tại $B$:
			$$AB = \sqrt{OA^2 - OB^2} = \sqrt{(3R)^2 - R^2} = 2R\sqrt{2}.$$
			Đoạn cao $BH$ hạ từ $B$ xuống cạnh huyền $OA$:
			$$BH = \dfrac{AB \cdot OB}{OA} = \dfrac{2R\sqrt{2} \cdot R}{3R} = \dfrac{2R\sqrt{2}}{3}.$$
			Vì $H$ là trung điểm của dây cung $BC$ nên:
			$$BC = 2BH = 2 \cdot \dfrac{2R\sqrt{2}}{3} = \dfrac{4R\sqrt{2}}{3}.$$
			Diện tích tam giác $ABO$ là:
			$$S_{ABO} = \dfrac{1}{2} AB \cdot OB = \dfrac{1}{2} \cdot 2R\sqrt{2} \cdot R = R^2\sqrt{2}.$$
			Diện tích tứ giác $ABOC$ là:
			$$S_{ABOC} = 2S_{ABO} = 2R^2\sqrt{2}.$$
			Vậy $BC = \dfrac{4R\sqrt{2}}{3}$ và $S_{ABOC} = 2R^2\sqrt{2}$.
		\end{enumerate}
	}
\end{bt}

% ==============================================================================
% CỤM 9 - ĐỀ THAM KHẢO 1 (Đề số 25)
% ==============================================================================

\begin{bt}%[9D6H18-2]
% nguon: Đề tham khảo Tuyển sinh 10 năm học 2026-2027 Sở GD&ĐT TP.HCM - Cụm 9 Đề 1, Bài 1
	Cho hàm số $y = \dfrac{x^2}{2}$ có đồ thị là $(P)$.
	\begin{enumerate}[a)]
		\item Vẽ đồ thị $(P)$ trên mặt phẳng tọa độ $Oxy$.
		\item Tìm các điểm thuộc $(P)$ khác gốc tọa độ và có hoành độ gấp đôi tung độ.
	\end{enumerate}
	\loigiai{
		\begin{enumerate}[a)]
			\item Bảng giá trị của hàm số $y = \dfrac{x^2}{2}$:
			\begin{center}
				\begin{tabular}{|c|c|c|c|c|c|}
					\hline
					$x$ & $-4$ & $-2$ & $0$ & $2$ & $4$ \\
					\hline
					$y = \dfrac{x^2}{2}$ & $8$ & $2$ & $0$ & $2$ & $8$ \\
					\hline
				\end{tabular}
			\end{center}
			Đồ thị $(P)$ là một parabol đi qua gốc tọa độ $O(0;0)$, nằm phía trên trục hoành và nhận trục tung $Oy$ làm trục đối xứng.
			\begin{center}
				\begin{tikzpicture}[>=stealth, scale=0.7]
					\draw[->] (-5,0) -- (5,0) node[below] {$x$};
					\draw[->] (0,-1) -- (0,9) node[left] {$y$};
					\fill (0,0) circle (1.5pt) node[below right] {$O$};
					\draw[smooth, samples=100, domain=-4.1:4.1, thick, color=blue] plot(\x, {1/2*(\x)^2});
					\draw[dashed] (-4,0) node[below] {$-4$} -- (-4,8) -- (0,8) node[left] {$8$} -- (4,8) -- (4,0) node[below] {$4$};
					\draw[dashed] (-2,0) node[below] {$-2$} -- (-2,2) -- (0,2) node[left] {$2$} -- (2,2) -- (2,0) node[below] {$2$};
					\fill (-4,8) circle (1.5pt) (4,8) circle (1.5pt) (-2,2) circle (1.5pt) (2,2) circle (1.5pt);
				\end{tikzpicture}
			\end{center}
			\item Vì điểm cần tìm có hoành độ gấp đôi tung độ nên $x = 2y \iff y = \dfrac{1}{2}x$.
			
			Phương trình hoành độ giao điểm:
			$$\dfrac{x^2}{2} = \dfrac{1}{2}x \iff x^2 - x = 0 \iff x(x - 1) = 0 \iff \left[\begin{aligned} &x = 0 \\ &x = 1. \end{aligned}\right.$$
			Vì điểm cần tìm khác gốc tọa độ $O$ nên $x \neq 0$, suy ra $x = 1$.
			
			Với $x = 1 \implies y = \dfrac{1^2}{2} = 0{,}5$.
			
			Vậy tọa độ điểm cần tìm là $(1; 0{,}5)$.
		\end{enumerate}
	}
\end{bt}

\begin{bt}%[9D6H20-2]
% nguon: Đề tham khảo Tuyển sinh 10 năm học 2026-2027 Sở GD&ĐT TP.HCM - Cụm 9 Đề 1, Bài 2
	Cho phương trình $x^2 - 2x - 3 = 0$.
	\begin{enumerate}[a)]
		\item Chứng minh phương trình luôn có hai nghiệm phân biệt $x_1, x_2$.
		\item Không giải phương trình, hãy tính giá trị của biểu thức:
		$$Q = x_1 x_2 - 4x_2(x_1 - x_2) + 4x_1^2.$$
	\end{enumerate}
	\loigiai{
		\begin{enumerate}[a)]
			\item Phương trình $x^2 - 2x - 3 = 0$ có các hệ số $a = 1, b = -2, c = -3$.
			
			Vì $a \cdot c = 1 \cdot (-3) = -3 < 0$ nên phương trình luôn có hai nghiệm phân biệt $x_1, x_2$.
			\item Theo định lí Viète, ta có:
			$$\begin{cases} x_1 + x_2 = 2 \\ x_1 x_2 = -3. \end{cases}$$
			Ta có:
			$$x_1^2 + x_2^2 = (x_1 + x_2)^2 - 2x_1 x_2 = 2^2 - 2 \cdot (-3) = 4 + 6 = 10.$$
			Biến đổi biểu thức $Q$:
			$$Q = x_1 x_2 - 4x_1 x_2 + 4x_2^2 + 4x_1^2 = 4(x_1^2 + x_2^2) - 3x_1 x_2.$$
			Thay các giá trị vào biểu thức $Q$, ta được:
			$$Q = 4 \cdot 10 - 3 \cdot (-3) = 40 + 9 = 49.$$
			Vậy $Q = 49$.
		\end{enumerate}
	}
\end{bt}

\begin{bt}%[9D8H26-3]
% nguon: Đề tham khảo Tuyển sinh 10 năm học 2026-2027 Sở GD&ĐT TP.HCM - Cụm 9 Đề 1, Bài 3
	Khảo sát thời gian tiếp xúc với các thiết bị điện tử trung bình trong một ngày của học sinh một trường THCS được cho như sau: $0{,}5$ giờ có $18$ học sinh; $1$ giờ có $46$ học sinh; $2$ giờ có $68$ học sinh; $3$ giờ có $51$ học sinh; $4$ giờ có $17$ học sinh.
	\begin{enumerate}[a)]
		\item Tính tổng số học sinh tham gia cuộc khảo sát.
		\item Chọn ngẫu nhiên một học sinh. Tính xác suất của các biến cố sau:
		\begin{itemize}
			\item Biến cố $A$: ``Học sinh được chọn có thời gian tiếp xúc thiết bị điện tử trung bình trên $3$ giờ một ngày''.
			\item Biến cố $B$: ``Học sinh được chọn có thời gian tiếp xúc thiết bị điện tử theo đúng khuyến nghị (không quá $2$ giờ mỗi ngày)''.
		\end{itemize}
	\end{enumerate}
	\loigiai{
		\begin{enumerate}[a)]
			\item Tổng số học sinh tham gia khảo sát là:
			$$n(\Omega) = 18 + 46 + 68 + 51 + 17 = 200\text{ (học sinh)}.$$
			\item Xác suất của từng biến cố:
			\begin{itemize}
				\item Học sinh có thời gian tiếp xúc trên $3$ giờ một ngày là các em tiếp xúc $4$ giờ/ngày ($17$ học sinh).
				
				Số kết quả thuận lợi cho biến cố $A$ là $n(A) = 17$.
				
				Xác suất của biến cố $A$ là:
				$$P(A) = \dfrac{n(A)}{n(\Omega)} = \dfrac{17}{200} = 0{,}085.$$
				\item Học sinh tiếp xúc không quá $2$ giờ mỗi ngày gồm các mức $0{,}5$ giờ, $1$ giờ và $2$ giờ.
				
				Số kết quả thuận lợi cho biến cố $B$ là:
				$$n(B) = 18 + 46 + 68 = 132\text{ (học sinh)}.$$
				Xác suất của biến cố $B$ là:
				$$P(B) = \dfrac{n(B)}{n(\Omega)} = \dfrac{132}{200} = \dfrac{33}{50} = 0{,}66.$$
			\end{itemize}
		\end{enumerate}
	}
\end{bt}

\begin{bt}%[9D6V21-3]
% nguon: Đề tham khảo Tuyển sinh 10 năm học 2026-2027 Sở GD&ĐT TP.HCM - Cụm 9 Đề 1, Bài 4
	Để thiết kế mô hình đồ chơi con vịt, bạn Nam vẽ hình vuông $ABCD$ cạnh dài $30\text{ cm}$. Phần thân con vịt gồm hình quạt tròn $ABD$ (cung $BD$ tâm $A$, góc ở tâm $90^\circ$) khoét đi nửa hình tròn đường kính $AE = 2x\text{ (cm)}$ ($0 < x < 15$). Phần đầu vịt là nửa hình tròn đường kính $EB$ và phần sau vịt là nửa hình tròn đường kính $AD$. Lấy $\pi \approx 3{,}14$.
	\begin{enumerate}[a)]
		\item Viết biểu thức tính diện tích phần thân mô hình con vịt theo $x$.
		\item Tính diện tích toàn bộ mô hình con vịt biết diện tích phần thân là $549{,}5\text{ cm}^2$.
	\end{enumerate}
	\loigiai{
		\begin{enumerate}[a)]
			\item Diện tích hình quạt tròn $ABD$ tâm $A$, bán kính $R = 30\text{ cm}$ là:
			$$S_{\text{quạt}} = \dfrac{\pi R^2 \cdot 90^\circ}{360^\circ} = \dfrac{\pi \cdot 30^2}{4} = 225\pi \approx 225 \cdot 3{,}14 = 706{,}5\text{ (cm}^2).$$
			Bán kính của nửa hình tròn đường kính $AE$ là:
			$$r_1 = \dfrac{2x}{2} = x\text{ (cm)}.$$
			Diện tích nửa hình tròn đường kính $AE$ là:
			$$S_1 = \dfrac{1}{2}\pi x^2 \approx 1{,}57x^2\text{ (cm}^2).$$
			Biểu thức tính diện tích phần thân mô hình con vịt là:
			$$S_{\text{thân}} = S_{\text{quạt}} - S_1 = 225\pi - \dfrac{1}{2}\pi x^2 \approx 706{,}5 - 1{,}57x^2\text{ (cm}^2).$$
			\item Theo đề bài, diện tích phần thân là $549{,}5\text{ cm}^2$ nên ta có:
			$$706{,}5 - 1{,}57x^2 = 549{,}5 \iff 1{,}57x^2 = 157 \iff x^2 = 100 \implies x = 10\text{ (cm) (do } x > 0).$$
			Khi $x = 10\text{ cm}$, độ dài đoạn $AE = 2x = 20\text{ cm}$.
			
			Đường kính phần đầu vịt là:
			$$EB = AB - AE = 30 - 20 = 10\text{ (cm)}.$$
			Bán kính phần đầu vịt là $r_2 = 5\text{ cm}$. Diện tích phần đầu vịt là:
			$$S_{\text{đầu}} = \dfrac{1}{2}\pi \cdot 5^2 = 12{,}5\pi\text{ (cm}^2).$$
			Đường kính phần sau vịt là $AD = 30\text{ cm}$, bán kính là $r_3 = 15\text{ cm}$. Diện tích phần sau vịt là:
			$$S_{\text{sau}} = \dfrac{1}{2}\pi \cdot 15^2 = 112{,}5\pi\text{ (cm}^2).$$
			Tổng diện tích phần đầu và phần sau vịt là:
			$$S_{\text{đầu}} + S_{\text{sau}} = 12{,}5\pi + 112{,}5\pi = 125\pi \approx 125 \cdot 3{,}14 = 392{,}5\text{ (cm}^2).$$
			Diện tích toàn bộ mô hình con vịt là:
			$$S = S_{\text{thân}} + S_{\text{đầu}} + S_{\text{sau}} = 549{,}5 + 392{,}5 = 942\text{ (cm}^2).$$
			Vậy diện tích toàn bộ mô hình con vịt là $942\text{ cm}^2$.
		\end{enumerate}
	}
\end{bt}

\begin{bt}%[9H0V31-3]
% nguon: Đề tham khảo Tuyển sinh 10 năm học 2026-2027 Sở GD&ĐT TP.HCM - Cụm 9 Đề 1, Bài 5
	Một quả cầu pha lê hình cầu có diện tích mặt cầu là $196\pi\text{ cm}^2$.
	\begin{enumerate}[a)]
		\item Tính bán kính của quả cầu pha lê.
		\item Người ta đặt quả cầu pha lê chìm hoàn toàn vào một bể cá hình hộp chữ nhật có kích thước đáy $25\text{ cm} \times 40\text{ cm}$ và chiều cao $50\text{ cm}$. Bể cá lúc đầu đang chứa sẵn $40$ lít nước. Tính khoảng cách từ mép nước đến miệng thành bể sau khi thả quả cầu (làm tròn kết quả đến hàng phần mười của $\text{cm}$).
	\end{enumerate}
	\loigiai{
		\begin{enumerate}[a)]
			\item Diện tích mặt cầu là $S = 4\pi R^2$.
			
			Theo đề bài ta có:
			$$4\pi R^2 = 196\pi \iff R^2 = 49 \implies R = 7\text{ (cm)}.$$
			Vậy bán kính của quả cầu pha lê là $7\text{ cm}$.
			\item Thể tích của quả cầu pha lê là:
			$$V_{\text{cầu}} = \dfrac{4}{3}\pi R^3 = \dfrac{4}{3}\pi \cdot 7^3 = \dfrac{1372}{3}\pi \approx \dfrac{1372 \cdot 3{,}14}{3} \approx 1435{,}97\text{ (cm}^3).$$
			Lượng nước ban đầu trong bể là:
			$$V_{\text{nước}} = 40\text{ lít} = 40\,000\text{ (cm}^3).$$
			Tổng thể tích của nước và quả cầu trong bể là:
			$$V = V_{\text{nước}} + V_{\text{cầu}} = 40\,000 + 1435{,}97 = 41435{,}97\text{ (cm}^3).$$
			Diện tích đáy của bể cá hình hộp chữ nhật là:
			$$S_{\text{đáy}} = 25 \cdot 40 = 1000\text{ (cm}^2).$$
			Chiều cao của mực nước sau khi thả quả cầu là:
			$$h_{\text{nước}} = \dfrac{V}{S_{\text{đáy}}} = \dfrac{41435{,}97}{1000} \approx 41{,}44\text{ (cm)}.$$
			Khoảng cách từ mép nước đến miệng thành bể lúc này là:
			$$d = 50 - 41{,}44 = 8{,}56 \approx 8{,}6\text{ (cm)}.$$
			Vậy khoảng cách từ mép nước đến miệng bể là khoảng $8{,}6\text{ cm}$.
		\end{enumerate}
	}
\end{bt}

\begin{bt}%[9D1H3-4]
% nguon: Đề tham khảo Tuyển sinh 10 năm học 2026-2027 Sở GD&ĐT TP.HCM - Cụm 9 Đề 1, Bài 6
	Hai xe buýt chạy trên một tuyến đường khép kín $ABCDA$ liên tục không nghỉ với vận tốc không đổi. Biết $AB = 26\text{ km}$, $BC = 19\text{ km}$, $CD = 23\text{ km}$ và $DA = 32\text{ km}$. Nếu cùng xuất phát từ $A$ đi ngược chiều nhau thì sau $1{,}2$ giờ hai xe gặp nhau; nếu đi cùng chiều thì xe thứ hai sẽ vượt xe thứ nhất sau $6$ giờ. Tính thời gian ít nhất để mỗi xe đi từ $A$ đến $C$.
	\loigiai{
		Tổng chiều dài lộ trình một vòng khép kín là:
		$$S = AB + BC + CD + DA = 26 + 19 + 23 + 32 = 100\text{ (km)}.$$
		Gọi $v_1, v_2$ (km/h) lần lượt là vận tốc của xe thứ nhất và xe thứ hai ($v_2 > v_1 > 0$).
		\begin{itemize}
			\item Khi hai xe đi ngược chiều, sau $1{,}2$ giờ hai xe gặp nhau nên tổng quãng đường hai xe đi được đúng bằng $1$ vòng lộ trình:
			$$1{,}2(v_1 + v_2) = 100 \iff v_1 + v_2 = \dfrac{250}{3}. \quad (1)$$
			\item Khi hai xe đi cùng chiều, sau $6$ giờ xe thứ hai vượt xe thứ nhất nên hiệu quãng đường hai xe đi được bằng $1$ vòng lộ trình:
			$$6(v_2 - v_1) = 100 \iff v_2 - v_1 = \dfrac{50}{3}. \quad (2)$$
		\end{itemize}
		Từ $(1)$ và $(2)$, ta có hệ phương trình:
		$$\begin{cases} v_1 + v_2 = \dfrac{250}{3} \\ -v_1 + v_2 = \dfrac{50}{3} \end{cases} \iff \begin{cases} 2v_2 = 100 \\ v_1 = v_2 - \dfrac{50}{3} \end{cases} \iff \begin{cases} v_2 = 50\text{ (km/h)} \\ v_1 = \dfrac{100}{3}\text{ (km/h)}. \end{cases}$$
		Để đi từ $A$ đến $C$, có hai hướng đi:
		- Hướng qua $B$: $s_1 = AB + BC = 26 + 19 = 45\text{ km}$.
		- Hướng qua $D$: $s_2 = AD + DC = 32 + 23 = 55\text{ km}$.
		
		Quãng đường ngắn nhất từ $A$ đến $C$ là $45\text{ km}$.
		
		Thời gian ít nhất để xe thứ hai đi từ $A$ đến $C$ là:
		$$t_2 = \dfrac{45}{50} = 0{,}9\text{ (giờ)} = 54\text{ phút}.$$
		Thời gian ít nhất để xe thứ nhất đi từ $A$ đến $C$ là:
		$$t_1 = \dfrac{45}{\dfrac{100}{3}} = \dfrac{135}{100} = 1{,}35\text{ (giờ)} = 1\text{ giờ } 21\text{ phút}.$$
	}
\end{bt}

\begin{bt}%[9H9C29-3]
% nguon: Đề tham khảo Tuyển sinh 10 năm học 2026-2027 Sở GD&ĐT TP.HCM - Cụm 9 Đề 1, Bài 7
	Cho nửa đường tròn $(O; R)$ có đường kính $AB$. Trên nửa đường tròn lấy điểm $E$ tùy ý, tia phân giác của góc $\widehat{EBA}$ cắt nửa đường tròn tại $D$. Gọi $H$ là giao điểm của $AE$ và $BD$. Kẻ $HF \perp AB$ tại $F$.
	\begin{enumerate}[a)]
		\item Chứng minh $\widehat{BDA} = 90^\circ$ và $\widehat{HFD} = \widehat{HAD}$.
		\item Gọi $C$ là giao điểm của $BE$ và $AD$. Chứng minh ba điểm $C, H, F$ thẳng hàng và $DH \cdot DB = DF^2 = \dfrac{AC^2}{4}$.
		\item Tính các góc của tam giác $DEF$ và chu vi của tam giác $DEF$ theo $R$ trong trường hợp $\widehat{ABC} = 30^\circ$.
	\end{enumerate}
	\loigiai{
		\begin{center}
			\begin{tikzpicture}[>=stealth, scale=1]
				\coordinate (A) at (-3,0);
				\coordinate (B) at (3,0);
				\coordinate (O) at (0,0);
				\def\r{3}
				\draw (A) -- (B);
				\draw (A) arc (180:0:\r);
				% E trên nửa đường tròn
				\coordinate (E) at (-1.5, 2.598);
				\draw (A) -- (E);
				\draw (B) -- (E);
				% BD là phân giác góc EBA
				% Góc EBA = 60 độ => phân giác 30 độ
				\coordinate (D) at (60:\r);
				\draw (B) -- (D);
				\draw (A) -- (D);
				\coordinate (H) at (intersection of A--E and B--D);
				\coordinate (F) at ($(A)!(H)!(B)$);
				\draw (H) -- (F);
				\coordinate (C) at (intersection of B--E and A--D);
				\draw (E) -- (C);
				\draw (D) -- (C);
				\draw (C) -- (H);
				\draw (D) -- (F);
				\draw (D) -- (E);
				\draw (E) -- (F);
				\node at (A) [below left] {$A$};
				\node at (B) [below right] {$B$};
				\node at (O) [below] {$O$};
				\node at (E) [above left] {$E$};
				\node at (D) [above right] {$D$};
				\node at (H) [below left] {$H$};
				\node at (F) [below] {$F$};
				\node at (C) [above] {$C$};
				\fill (A) circle (1.5pt) (B) circle (1.5pt) (O) circle (1.5pt) (E) circle (1.5pt) (D) circle (1.5pt) (H) circle (1.5pt) (F) circle (1.5pt) (C) circle (1.5pt);
			\end{tikzpicture}
		\end{center}
		\begin{enumerate}[a)]
			\item Vì $\widehat{BDA}$ là góc nội tiếp chắn nửa đường tròn đường kính $AB$ nên $\widehat{BDA} = 90^\circ$.
			
			Do đó $AD \perp BD \implies \widehat{HDA} = 90^\circ$.
			
			Vì $HF \perp AB$ tại $F$ nên $\widehat{HFA} = 90^\circ$.
			
			Tứ giác $AFHD$ có $\widehat{HDA} = \widehat{HFA} = 90^\circ$, do đó hai đỉnh kề nhau $D$ và $F$ cùng nhìn đoạn $AH$ dưới một góc vuông.
			
			Suy ra tứ giác $AFHD$ nội tiếp đường tròn đường kính $AH$.
			
			Do đó $\widehat{HFD} = \widehat{HAD}$ (hai góc nội tiếp cùng chắn cung $HD$).
			\item Vì $\widehat{AEB} = 90^\circ$ (góc nội tiếp chắn nửa đường tròn) nên $AE \perp BC$.
			
			Mặt khác $AD \perp BC$ tại $D$ nên trong tam giác $ABC$, $AE$ và $BD$ là hai đường cao cắt nhau tại $H$.
			
			Do đó $H$ là trực tâm của tam giác $ABC$, suy ra $CH \perp AB$.
			
			Mà $HF \perp AB$ tại $F$ nên ba điểm $C, H, F$ thẳng hàng.
			
			Vì $BD$ là phân giác của góc $\widehat{EBA}$ nên trong tam giác $ABC$, đoạn $BD$ vừa là đường cao vừa là đường phân giác.
			
			Suy ra tam giác $ABC$ cân tại $B$, do đó $BD$ đồng thời là đường trung tuyến của tam giác $ABC$.
			
			Suy ra $D$ là trung điểm của cạnh $AC$, do đó $AD = DC = \dfrac{AC}{2}$.
			
			Trong tam giác vuông $AFC$ tại $F$, có $D$ là trung điểm của cạnh huyền $AC$, nên $FD$ là đường trung tuyến ứng với cạnh huyền:
			$$DF = \dfrac{AC}{2} \implies DF^2 = \dfrac{AC^2}{4}.$$
			Xét $\Delta ADH$ và $\Delta BDA$ có:
			$$\widehat{ADH} = \widehat{BDA} = 90^\circ,\quad \widehat{DAH} = \widehat{DBA}\text{ (cùng phụ với }\widehat{DAB}).$$
			Do đó $\Delta ADH \backsim \Delta BDA\text{ (g-g)}$, suy ra:
			$$\dfrac{DH}{AD} = \dfrac{AD}{BD} \implies DH \cdot DB = AD^2 = \left(\dfrac{AC}{2}\right)^2 = \dfrac{AC^2}{4}.$$
			Vậy $DH \cdot DB = DF^2 = \dfrac{AC^2}{4}$.
			\item Khi $\widehat{ABC} = 30^\circ$:
			
			Vì tam giác $ABC$ cân tại $B$ nên $\widehat{BAC} = \widehat{BCA} = \dfrac{180^\circ - 30^\circ}{2} = 75^\circ$.
			
			Vì $BD$ là phân giác của $\widehat{ABC}$ nên $\widehat{ABD} = \dfrac{30^\circ}{2} = 15^\circ$.
			
			Trong tam giác vuông $BDA$ tại $D$:
			$$AD = AB \cdot \sin 15^\circ = 2R \cdot \sin 15^\circ.$$
			Ta có $D$ là điểm chính giữa cung $AE$ nên $DE = DA = 2R\sin 15^\circ$.
			
			Tam giác $DEF$ có $FD = \dfrac{AC}{2} = AD$ và tứ giác $DFEH$ nội tiếp, các góc của tam giác $DEF$ lần lượt là $60^\circ, 60^\circ, 60^\circ$ (tam giác $DEF$ đều).
			
			Chu vi tam giác $DEF$ là:
			$$C = 3 \cdot DE = 3 \cdot 2R\sin 15^\circ = 6R\sin 15^\circ = 6R \cdot \dfrac{\sqrt{6} - \sqrt{2}}{4} = \dfrac{3(\sqrt{6} - \sqrt{2})R}{2}.$$
		\end{enumerate}
	}
\end{bt}
